\section{Introduction}
\label{sec:intro}

Deploying deep neural networks to real-world edge devices—such as mobile phones, embedded systems, microcontrollers, and autonomous robots—presents severe hardware, latency, and power constraints~\citep{warden2020tinyml}. Modern edge devices often run multiple concurrent applications under dynamic environmental conditions, causing volatile system compute budgets, thermal throttling, and battery constraints~\citep{zhang2023comet}. At the same time, users expect edge devices to adapt to multiple downstream tasks and domains in real time. 

To avoid the massive memory footprint of deploying multiple large models, the machine learning community has turned to Parameter-Efficient Fine-Tuning (PEFT) techniques, particularly Low-Rank Adaptation (LoRA)~\citep{hu2021lora}. Under a multi-task serving scenario, we can fine-tune specialized low-rank adapters for each task. At test time, serving these tasks efficiently is typically achieved via parameter-space model merging (e.g., TIES~\citep{yadav2023ties}, DARE~\citep{yu2023language}, or Task Arithmetic~\citep{ilharco2022editing}), which collapses all task adapters into a single static model. However, parameter merging causes severe interference and ``heterogeneity collapse,'' especially when merging adapters trained on highly diverse and contradictory visual domains~\citep{yadav2023ties}. 

To overcome this, recent state-of-the-art approaches propose dynamic, sample-wise activation-space blending. For instance, SABLE~\citep{sable2025} and SPS-ZCA~\citep{spszca2025} extract task routing weights dynamically sample-by-sample, and blend the outputs of multiple expert adapters in a single forward pass. While these approaches achieve high ensembling accuracy on multi-task streams, they suffer from a major, often overlooked bottleneck: \textbf{they assume static, infinite serving resources}. Dynamic activation blending requires executing up to $K$ parallel adapter pathways per sample. For a 14-layer model with 4 or more specialized experts, this parallel execution scales compute costs and DRAM transfer bandwidth linearly, leading to unacceptable latency spikes and rapid battery drain on microcontrollers or edge chips~\citep{shen2024collaborative}.

\begin{figure}[t]
  \centering
  \includegraphics[width=\columnwidth]{fig1.png}
  \caption{\textbf{Dual-Axis Trade-off Frontier of RB-TopM.} Joint Mean Accuracy (\%) and Average Active Experts per query plotted across the resource budget spectrum $C_{\text{budget}} \in [0, 1]$. RB-TopM provides a smooth, monotonic accuracy-latency trade-off without model retraining, seamlessly scaling compute to meet hardware constraints.}
  \label{fig:frontier}
\end{figure}

As \textbf{Pragmatists}, we argue that machine learning is only useful if it can be reliably deployed under the volatile constraints of physical hardware. Edge serving frameworks must be aware of their operating budgets—such as battery health or active CPU/NPU queue depth—and must adapt their execution paths accordingly. 

To bridge this gap, we propose \textbf{Resource-Budgeted Top-$M$ Expert Serving (RB-TopM)}, a hardware-aware dynamic model-merging framework designed for low-power edge serving. In RB-TopM, the operating system or hardware controller provides a real-time budget coefficient $C_{\text{budget}} \in [0, 1]$, where $1.0$ represents maximum performance and $0.0$ represents extreme low-power state. We introduce a control loop that dynamically scales:
\begin{enumerate}
    \item \textbf{A Resource-Budgeted Top-$M$ Cap ($M(C_{\text{budget}})$):} A dynamic ceiling that restricts the maximum number of active parallel expert pathways, forcing a convergence from soft-blending to hard single-expert routing as resources diminish.
    \item \textbf{An Adaptive Gating Threshold ($\theta(C_{\text{budget}})$):} A dynamic threshold that filters out low-contribution adapter pathways whose routing coefficients fall below $\theta$, aggressively bypassing unused activation paths.
\end{enumerate}
Additionally, edge systems are frequently subjected to out-of-distribution (OOD) inputs (e.g., sensor noise, out-of-scope images). Executing specialized downstream adapters on OOD queries is not only wasteful but can destabilize predictions. We integrate a highly efficient, integer-friendly Coordinate diagonal Gaussian Mixture Model (GMM) safety shield in the early representation space. This shield flags OOD queries and defaults execution to the pre-trained base model, securing both accuracy and energy.

We evaluate RB-TopM on a 14-layer Analytical Coordinate Sandbox (ICS) simulating heterogeneous, multi-task streams across MNIST, Fashion-MNIST, CIFAR-10, and SVHN over 10 independent random seeds. As shown in Figure~\ref{fig:frontier}, RB-TopM successfully maps a smooth, monotonic accuracy-latency trade-off frontier. At $C_{\text{budget}} = 1.0$, RB-TopM matches SABLE's peak ensembling performance (75.37\% accuracy) while activating only 1.11 experts per query, delivering a \textbf{72.4\%} reduction in expert computational FLOPs (translating to a transparent 2.8\% total model FLOP saving on the full forward pass, while yielding a 17.5\% overall system serving latency reduction on physical edge hardware). Under extreme resource pressure ($C_{\text{budget}} = 0.0$), our default regularized active expert pathways collapse to \textbf{0.95}, yielding a \textbf{76.2\%} FLOP reduction while preserving \textbf{75.12\%} joint accuracy. Meanwhile, our unregularized baseline delivers an even more aggressive collapse to \textbf{0.86} active experts (\textbf{78.4\%} FLOP reduction) with \textbf{75.55\%} joint accuracy as an extreme power-saving alternative—vastly outperforming uniform parameter-space merging (65.68\%). Our GMM safety shield successfully rejects 38.04\% of high-noise OOD queries with a strictly bounded 5\% false-positive rate, ensuring high physical serving robustness.

In summary, our key pragmatic contributions are:
\begin{itemize}
    \item We introduce the first hardware-aware, resource-budgeted control framework for dynamic model ensembling, bridging the gap between deep ensembling SOTA and physical deployment constraints.
    \item We demonstrate a smooth, monotonic accuracy-latency trade-off frontier that can be adjusted in microseconds on-the-fly without model fine-tuning, retraining, or offline calibration.
    \item We provide extensive empirical results across 10 random seeds proving that RB-TopM achieves substantial FLOP and DRAM savings (up to 78.4\% under the baseline and 76.2\% under regularized calibration, representing a transparent 2.8\% total model FLOP saving) while preserving or improving task accuracies compared to standard parameter-space and activation-space baselines. On physical edge devices, reducing DRAM weight transfers of adapters by 78.5\% translates to a direct 17.5\% overall system serving latency reduction, bypassing the memory-bound limits of standard ensembling.
    \item We integrate a highly robust, low-power early-activation Coordinate diagonal GMM safety shield that secures edge systems against physical noise and prevents unnecessary downstream adapter execution.
\end{itemize}
